\section{핵심 가설}
\label{sec:hyp}

\begin{itemize}
  \item \textbf{H-J (짐 덜기).} 계획·판단·명령을 상위 VLM이 모두 맡고 VLA는 조이스틱(명령 기준 미세 조종)만 하면, 같은 VLA가 혼자 할 때보다 낫다. 검증은 같은 VLA를 단독 모드(스스로 계획)와 조이스틱 모드로 같은 과제·시드에서 돌려 비교한다(\cref{sec:eval}). 단독 모드는 비교용으로만 둔다.
  \item \textbf{H-L (학습 가능한 상위).} 학습 가능한 크기(35B-A3B, LoRA)의 VLM을 시뮬 참값(특권 정보로 계산)과 공개 로봇 데이터의 점 라벨로 특화하면, 로봇을 조종하는 상위로 쓸 만하다.
  \item \textbf{H-C (연결 손실).} 잘 설계한 연결은 두 모듈의 성능을 거의 깎지 않는다. 연결 손실(\cref{sec:closs}) $\le$ 5\,\%p, 연결 탓 실패 $\le$ 전체 실패의 10\,\%를 합격 기준으로 둔다.
\end{itemize}
H-L은 지금 근거가 가장 많고(\cref{sec:evidence}), H-J와 H-C는 결합 실험(\cref{sec:coupling})이 직접 겨누는 가설이다.
